% titel: Blatt 2, A5: Summenformel aus Massenanteilen und Molmasse
% Teil II: der durchgerechnete Weg.

\subsection*{Gegeben und gesucht}
Zusammensetzung: $47{,}3\%$ C, $2{,}54\%$ H, $50{,}0\%$ Cl; experimentelle
Molmasse $M = \qty{355}{g/mol}$. Gesucht: Summenformel.

\subsection*{Masse jedes Elements in einem mol}
\begin{align*}
  m_{\mathrm{C}} &= 0{,}473 \cdot 355 = 167{,}9\ \mathrm{g},\\
  m_{\mathrm{H}} &= 0{,}0254 \cdot 355 = 9{,}02\ \mathrm{g},\\
  m_{\mathrm{Cl}} &= 0{,}500 \cdot 355 = 177{,}5\ \mathrm{g}.
\end{align*}

\subsection*{Stoffmengen}
\begin{align*}
  n_{\mathrm{C}} &= \frac{167{,}9}{12{,}011} = 13{,}98\ \mathrm{mol},\\
  n_{\mathrm{H}} &= \frac{9{,}02}{1{,}008} = 8{,}95\ \mathrm{mol},\\
  n_{\mathrm{Cl}} &= \frac{177{,}5}{35{,}453} = 5{,}01\ \mathrm{mol}.
\end{align*}
Verhältnis durch $n_{\mathrm{Cl}}$ teilen: $13{,}98:8{,}95:5{,}01
\approx 14:9:5$.

\begin{ergebnis}
\begin{equation}
  \mathrm{C_{14}H_{9}Cl_{5}},\qquad
  M_{\mathrm{theor}} = 14\cdot12{,}011 + 9\cdot1{,}008
    + 5\cdot35{,}453 = \qty{354}{g/mol}.
\end{equation}
$M_{\mathrm{theor}} \approx 354\ \mathrm{g/mol}$ stimmt mit dem gemessenen Wert
$355\ \mathrm{g/mol}$ überein.
\end{ergebnis}

\begin{kontrolle}
Massenanteile zurückrechnen: $14\cdot12{,}011/354 = 47{,}5\%$ C,
$9\cdot1{,}008/354 = 2{,}6\%$ H und $5\cdot35{,}453/354 = 50{,}1\%$ Cl --- sehr
gut passend. Die theoretische Molmasse liegt in der Nahenordnung der
experimentell bestimmten ($354$ gegen $355$ g/mol); der Rest ist Messungauigkeit.
\end{kontrolle}

\begin{merksatz}
Aus $w_i$ und der Molmasse zuerst die Masse je Element im mol, dann die
Stoffmengen bilden, ins ganzzahlige Verhältnis bringen und mit der Molmasse
auf konsistenz prüfen. Empirische und Summenformel können sich um einen
ganzzahligen Faktor unterscheiden.
\end{merksatz}
